\documentclass[9pt,a4paper]{extarticle}

% ---------- THEME: Packages & Geometry ----------
\usepackage[left=0.25in,right=0.25in,top=0.18in,bottom=0.08in]{geometry}
\usepackage[T1]{fontenc}
\usepackage[table]{xcolor}
% Body font: keep exactly ONE of these lines uncommented.
% Carlito is very narrow; the alternatives are scaled so the page still fits on one page.
% Each scale is the largest that keeps this content on one page; re-check \SideLabel offsets after switching.
\usepackage[sfdefault,lf]{carlito} % 'lf' for lining figures
% \usepackage[default,scale=0.95]{lato}
% \usepackage[sfdefault,scale=0.92]{roboto}
% \usepackage[sfdefault,lining,scale=0.90]{FiraSans}
% \usepackage[default,scale=0.88]{opensans}
% \usepackage[sfdefault,scale=0.88]{noto-sans}
\usepackage{tcolorbox}
\usepackage{tabularx}
\usepackage{enumitem}
\usepackage{fontawesome5}
\usepackage{hyperref}
\usepackage{array}

% ---------- THEME: Fonts ----------
\newcommand{\HeadingFont}{\bfseries} 
\newcommand{\LinkFont}{\ttfamily} % cmtt has no bold weight

\newcommand{\NameFont}{\fontsize{18}{20}\selectfont\HeadingFont}
\newcommand{\BannerTitleFont}{\fontsize{11}{13}\selectfont\HeadingFont}
\newcommand{\EntryTitleFont}{\fontsize{9.5}{10.5}\selectfont\HeadingFont}
\newcommand{\BodyFontSize}{\fontsize{9}{9.5}\selectfont}
\newcommand{\LinkSize}{\fontsize{8}{9.5}\selectfont} % cmtt ships at 8pt, not 8.5pt

\newcommand{\CellStrut}{\rule[-0.6ex]{0pt}{2.4ex}}
\newcommand{\SideLabel}[2][0pt]{\raisebox{-#1}[0pt][0pt]{#2}}

% ---------- THEME: Colours ----------
% Accent colour: keep exactly ONE PrimaryColor line uncommented.
% All options keep white banner text readable (contrast >= 5.8:1).
\definecolor{PrimaryColor}{RGB}{47,84,150} % Navy
% \definecolor{PrimaryColor}{RGB}{30,95,140}  % Steel Blue
% \definecolor{PrimaryColor}{RGB}{60,80,110}  % Slate Blue
% \definecolor{PrimaryColor}{RGB}{0,105,120}  % Teal
% \definecolor{PrimaryColor}{RGB}{34,102,68}  % Forest Green
% \definecolor{PrimaryColor}{RGB}{88,50,140}  % Royal Purple
% \definecolor{PrimaryColor}{RGB}{128,28,48}  % Burgundy
% \definecolor{PrimaryColor}{RGB}{165,25,45}  % Crimson
% \definecolor{PrimaryColor}{RGB}{170,70,20}  % Burnt Orange
% \definecolor{PrimaryColor}{RGB}{55,60,70}   % Charcoal
\colorlet{SecondaryColor}{PrimaryColor!20} % row shading; higher (e.g. !25) = darker, lower (e.g. !10) = lighter
\definecolor{DarkGrey}{gray}{0.2}
\definecolor{MediumGrey}{gray}{0.45} 

\colorlet{BorderColor}{PrimaryColor}
\colorlet{BannerTextColor}{white}
\colorlet{LinkColor}{MediumGrey}

% ---------- THEME: Sizes & Spacing ----------
\newlength{\RuleWidth}
\setlength{\RuleWidth}{0.5pt}

\setlength{\tabcolsep}{3pt} % Reduced side padding
\renewcommand{\arraystretch}{0.95}\setlength{\extrarowheight}{1.5pt}
\arrayrulecolor{BorderColor}

\newlength{\LabelWidth}
\setlength{\LabelWidth}{2.8cm}

\newlength{\HeaderGap}
\setlength{\HeaderGap}{6pt}

\setlength{\parindent}{0pt}
\pagestyle{empty}

% Configure hyperref
\hypersetup{
    colorlinks=true,
    urlcolor=LinkColor,
    linkcolor=LinkColor,
    pdfborder={0 0 0},
    pdftitle={Resume -- Shubham Jadhav},
    pdfauthor={Shubham Jadhav}
}

\input{glyphtounicode}
\pdfgentounicode=1

% ---------- COMPONENTS ----------

% 1. Header (Aligned as a single strip)
\newcommand{\ResumeHeader}[6]{%
  \noindent
  \parbox[b]{0.65\linewidth}{%
    {\NameFont #1 \,|\, #2} \\[2pt]
    {\BodyFontSize\color{DarkGrey} #3}%
  }%
  \hfill
  \parbox[b]{0.35\linewidth}{%
    \raggedleft
    \BodyFontSize
    #4 \\
    #5 \\
    #6%
  }%
  \vspace{\HeaderGap}\par
}

% 2. Section Block
\newtcolorbox{ResumeSection}[1]{
  colback=white,
  colframe=BorderColor,
  coltitle=BannerTextColor,
  boxrule=\RuleWidth,
  left=0pt, right=0pt, top=0pt, bottom=0pt,
  boxsep=0pt,
  before=\vspace{-0.5pt}\par\noindent,
  after=\par\nointerlineskip,
  fonttitle=\BannerTitleFont\vphantom{gjpqy},
  halign title=center,
  title={#1},
  toptitle=3pt, bottomtitle=3pt,
  titlerule=0pt,
  sharp corners
}

% 3. Bullets
\newenvironment{Bullets}{%
  \begin{itemize}[leftmargin=*, nosep, itemsep=2.5pt, topsep=1pt, partopsep=0pt, parsep=0pt, label={\small\textbullet}]
}{%
  \end{itemize}
}

% 4. Link & Date Range
\newcommand{\Link}[2]{\href{#1}{\LinkSize\LinkFont #2}}
\newcommand{\DateRange}[2]{#1 -- #2}

\begin{document}
\BodyFontSize

% ================= HEADER =================
\ResumeHeader{Shubham Jadhav}{CS24S009}%
{Department of Computer Science \& Engineering, IIT Madras}%
{\textbf{\faGlobe \hspace{2pt} Portfolio:} \Link{https://shubham-jadhav-FiatLux.github.io}{shubham-jadhav-FiatLux.github.io}}%
{\textbf{\faGithub \hspace{2pt} GitHub:} \Link{https://github.com/shubham-jadhav-FiatLux}{github.com/shubham-jadhav-FiatLux}}%


% ================= EDUCATION =================
\begin{ResumeSection}{Education}
  \begin{tabularx}{\linewidth}{ >{\centering\arraybackslash}m{2.5cm} | X | X | >{\centering\arraybackslash}m{2.5cm} }
    \rowcolor{SecondaryColor}
    \CellStrut\textbf{Year} & \textbf{Degree / Certificate} & \textbf{Institute} & \textbf{CGPA / \%} \\ \hline
    2024 -- Present & M.S. by Research / CSE & Indian Institute of Technology Madras & 7.78 / 10 \\ \hline
    2017 -- 2021 & B.Tech / CSE & Savitribai Phule Pune University & 9.08 / 10 \\ \hline
    2017 & Class XII (Maharashtra State Board) & Vivekanand College Kolhapur & 82.63 \% \\ \hline
    2015 & Class X (Maharashtra State Board) & Gadhinglaj High School Gadhinglaj & 94.20 \%
  \end{tabularx}
\end{ResumeSection}

% ================= TECHNICAL SKILLS =================
\begin{ResumeSection}{Technical Skills}
  \begin{tabularx}{\linewidth}{ >{\columncolor{SecondaryColor}\bfseries}p{\LabelWidth} | X }
    \CellStrut Languages & C and C++, Python, SQL \\ \hline
    \CellStrut HPC & CUDA, OpenMP, OpenACC, MPI, HIP \\ \hline
    \CellStrut Rendering & OpenGL, WebGL, DirectX, Three.js, GLSL, Win32 SDK \\ \hline
    \CellStrut Tools & Linux/POSIX, LLVM, Docker, Ghidra, Git/GitHub, LaTeX, Slurm, LLM, Ollama, LangChain \\ \hline
    \CellStrut \SideLabel[1.5\baselineskip]{Course Work} & 
    \textbf{PG Courses:} GPU Programming, Parallel Scientific Computing, Secure Systems Engineering, Program Analysis, Advanced Data Structures and Algorithms, Advanced Computer Architecture \par 
    \textbf{UG Courses:} Operating Systems, Data Structures and Algorithms, Computer Architecture, Database Management Systems, Computer Networks, Compiler Design
  \end{tabularx}
\end{ResumeSection}

% ================= RESEARCH WORK =================
\begin{ResumeSection}{Research}
  \begin{tabularx}{\linewidth}{ >{\columncolor{SecondaryColor}\bfseries}p{\LabelWidth} | X }
    \rowcolor{SecondaryColor}
    \multicolumn{2}{m{\dimexpr\linewidth-2\tabcolsep\relax}}{%
      \CellStrut{\EntryTitleFont Capacitated Vehicle Routing Problem} \hfill \DateRange{Jul'24}{Present} \newline
      \textit{M.S. Thesis | Guide: Prof. Rupesh Nasre, Prof. Narayanaswamy N S}
    } \\ \hline
    \CellStrut \SideLabel[2\baselineskip]{Context \& Problem} & The Capacitated Vehicle Routing Problem (CVRP) is an \textbf{NP-hard combinatorial optimization problem} which asks for a set of \textbf{minimum-cost routes} that start and end at a single depot. A homogeneous fleet of vehicles must service a set of customers with known demands, visiting each exactly once without exceeding vehicle capacity constraints. CVRP underpins modern logistics and distribution planning, where marginal savings in traveled distance yield massive \textbf{reductions in cost and emissions}. \\ \hline
    \CellStrut \SideLabel[3.0\baselineskip]{\parbox{\LabelWidth}{Approach \& \\ Contributions}} & 
    \vspace{-1.5ex}
    \begin{Bullets}
      \item \textbf{Engineered a high-performance C++ solver} utilizing a novel geometry-aware construction, highly granular local search, and multi-threaded ruin-and-recreate metaheuristics.
      \item \textbf{Pioneered an Adaptive Spiral Recombination architecture} to dynamically decompose and merge sub-problems, enabling the engine to scale efficiently on massive datasets where traditional routing algorithms asymptotically fail.
      \item Discovered \textbf{new Best-Known Solutions (BKS) on million-scale instances}, achieving a \textbf{10× runtime reduction} compared to SOTA.
      \item \textbf{Innovated a Geometry-Aware Initialization} that outperforms existing construction methods in both \textbf{speed and scalability}.
    \end{Bullets}
    \vspace{-1ex}
  \end{tabularx}
\end{ResumeSection}

% ================= PROJECTS =================
\begin{ResumeSection}{Projects}
  \begin{tabularx}{\linewidth}{ X }


    \rowcolor{SecondaryColor}
    \CellStrut{\EntryTitleFont TaintGate: Information-Flow Control for LLM Agents {\Link{https://github.com/shubham-jadhav-FiatLux}{\textbf \faGithub}} } \textit{| Python, MCP, LangChain, AgentDojo} \hfill {\DateRange{Jan'26}{May'26}} \newline 
    \textit{Secure Systems Engineering + Program Analysis + LLMs}
    \\
    \vspace{-1.5ex}
    \begin{Bullets}
      \item Engineered a deterministic \textbf{reference monitor} and \textbf{MCP proxy} to secure LLM agents against \textbf{prompt injection}, replacing brittle probabilistic guardrails.
      \item Enforced strict \textbf{information-flow control} via a \textbf{Biba $\times$ Bell-LaPadula} lattice, utilizing session-level \textbf{taint propagation} and backward \textbf{program slicing} to dynamically generate source-to-sink explanations for blocked actions.
      \item Evaluated defense efficacy via the \textbf{AgentDojo} benchmark (629 vectors), plotting precision-utility frontiers to quantify mitigation rates across models.
    \end{Bullets}
    \vspace{-1ex} \\ \hline

    \rowcolor{SecondaryColor}
    \CellStrut{\EntryTitleFont MiniDLC: End-to-End Deep Learning Compiler {\Link{https://github.com/shubham-jadhav-FiatLux}{\textbf \faGithub}} } \textit{| C++, Python, LLVM, CUDA} \hfill {\DateRange{Aug'25}{Dec'25}} \newline
    \textit{GPU Programming + Program Analysis + Computer Architecture}
    \\
    \vspace{-1.5ex}
    \begin{Bullets}
      \item Engineered an end-to-end compiler lowering models from a custom \textbf{Python frontend} into native machine code via C++ \textbf{SSA Graph} and \textbf{Loop IRs}.
      \item Implemented \textbf{reverse-mode autodiff} as a graph pass, optimized cache locality via dependence-checked \textbf{operator fusion}, and engineered liveness-based \textbf{memory planning} using interference coloring.
      \item Generated vectorized CPU code via \textbf{LLVM ORC JIT} and tiled GPU kernels via \textbf{CUDA NVRTC}, successfully training an \textbf{MNIST MLP} and \textbf{Char-GPT}.
    \end{Bullets}
    \vspace{-1ex} \\ \hline



    \rowcolor{SecondaryColor}
    \CellStrut{\EntryTitleFont Valley of Peace: Interactive 3D Portfolio {\Link{https://github.com/shubham-jadhav-FiatLux/shubham-jadhav-FiatLux.github.io}{\textbf \faGithub}} } \textit{| Three.js, GLSL, Web Audio} \hfill {\DateRange{Sep'21}{May'24}} \newline
    {\textit{ \hspace{2pt} Portfolio:} \Link{https://shubham-jadhav-FiatLux.github.io}{shubham-jadhav-FiatLux.github.io}}
    \\
    \vspace{-1.5ex}
    \begin{Bullets}
      \item Engineered an explorable, fully \textbf{procedural} 3D world (\textbf{21k LoC}) with \textbf{zero model, texture or audio files}, all generated in code.
      \item Rendered \textbf{1.7M triangles in $\sim$200 draw calls} via \textbf{GPU instancing} \& mesh batching, with \textbf{adaptive quality} for low-end devices.
      \item Directed a \textbf{5-min cinematic walkthrough} mode using an autopilot and a camera director on arc-length \textbf{B\'ezier splines}.
    \end{Bullets}
    \vspace{-1ex} \\ \hline



    \rowcolor{SecondaryColor}
    \CellStrut{\EntryTitleFont Other Projects} \\
    \vspace{-1.5ex}
    \begin{Bullets}
        \item \textbf{Coursework Implementations  {\Link{https://github.com/shubham-jadhav-FiatLux/IIT-Madras_Coursework}{\textbf \faGithub}} :} cache, branch-predictor, Tomasulo out-of-order \& \textbf{MESI} coherence simulators; \textbf{big.LITTLE} scheduler;
  \textbf{LLVM} range-analysis, Andersen points-to \& array bounds-check passes; \textbf{CUDA} Bor\r{u}vka MST and CPU-vs-GPU sine-wave simulation; stack-canary bypass,
  \textbf{ret2libc/ROP}, format-string \& glibc heap exploits; \textbf{MPI+OpenMP} ray tracer and parallel Poisson solvers.

        \item \textbf{Real-Time Rendering {\Link{https://github.com/shubham-jadhav-FiatLux/Real-Time_Rendering}{\textcolor{MediumGrey}{\faGithub}}} :} built \textbf{260+} rendering programs and \textbf{15+} mini-projects across \textbf{6 platforms} --- \textbf{Windows} (Win32
  SDK), \textbf{Linux} (X11/GLX), \textbf{Android} (Java), \textbf{macOS} \& \textbf{iOS} (Objective-C++), \textbf{Web} (JavaScript) --- covering projections, 3D
  transforms, texturing, model loading and per-fragment Phong lighting in \textbf{OpenGL} fixed-function \& programmable pipelines, \textbf{OpenGL ES}, \textbf{WebGL}
  and \textbf{Direct3D 11}, with \textbf{GLSL}/\textbf{HLSL} shaders and \textbf{CUDA}/\textbf{OpenCL}--OpenGL interop.
  {\textbf{{\textcolor{MediumGrey!70}{\faYoutube}} \hspace{2pt} YouTube :} \Link{https://www.youtube.com/@shubham-jadhav-FiatLux}{www.youtube.com/@shubham-jadhav-FiatLux}}%
        
      \item \textbf{PENETRALIA} -- Co-engineered a real-time graphics rendering demo featuring a CUDA-FFT ocean solver, procedural terrain, atmospheric scattering, and an advanced post-processing pipeline (godrays, bloom, depth of field) guided by B\'ezier camera animation. {\textcolor{MediumGrey!70}{\faYoutube}} \Link{https://youtu.be/Ze5QbGQnguM}{https://youtu.be/Ze5QbGQnguM}
      \item \textbf{Asa He Chanchal Maan} -- Demo created using the OpenGL graphics library and C language. Win32 SDK is used for windowing purposes and data structures in C are used to render scenes. {\textcolor{MediumGrey!70}{\faYoutube}} \Link{https://youtu.be/4pB2rZeWkNU}{https://youtu.be/4pB2rZeWkNU}
    \end{Bullets}
    \vspace{-1ex}
  \end{tabularx}
\end{ResumeSection}

% ================= POSITION OF RESPONSIBILITY =================
\begin{ResumeSection}{Positions of Responsibility \& Other Activities}
  \begin{tabularx}{\linewidth}{ >{\columncolor{SecondaryColor}\bfseries}p{\LabelWidth} | X }
    \CellStrut \SideLabel[0.5\baselineskip]{MS Representative} & Elected representative for \textbf{MS by Research scholars} in CSE at IIT Madras; collaborated with CSE faculty through class committee meetings to present student feedback on course delivery, quizzes, and assessments, and discuss course structure and grading. \\ \hline
    \CellStrut \SideLabel[0.5\baselineskip]{Teaching Assistant} & Served as TA for \textbf{Systems} and \textbf{Algorithms} courses, including GPU Programming, Blockchain Technology, and Natural Language Processing; designed problem sets, assignments and exams, conducted tutorials and office hours, and mentored students. \\ \hline
    \CellStrut \SideLabel[0.5\baselineskip]{Event Coordinator} & Coordinated logistics and technical operations for major IIT Madras events, including the \textbf{ICPC India Online Round} (2025, 2026) and the CSE Degree Distribution Ceremony. \\ \hline
    \CellStrut \SideLabel[0.5\baselineskip]{NSS Head} & Headed the \textbf{National Service Scheme} unit during B.Tech; directed community initiatives including street plays, blood donation drives, check-dam construction, tree plantation drives, and a 7-day rural development camp in Varoti Bk., Pune. \\ \hline
    \CellStrut RSP Leader & Led the \textbf{Road Safety Patrol} platoon during high school, overseeing team drills and public traffic safety awareness drives. \\ \hline
  \end{tabularx}
\end{ResumeSection}

% ================= INTERESTS & HOBBIES =================
\begin{ResumeSection}{Interests \& Hobbies}
  \begin{tabularx}{\linewidth}{ X }
    \CellStrut \textbf{Arts:} Drawing, Painting, Calligraphy \hfill | \hfill \textbf{Music:} Tabla, Guitar \hfill | \hfill \textbf{Sports:} Table Tennis, Cricket, Volleyball, Carrom, Chess, Kho-Kho
  \end{tabularx}
\end{ResumeSection}

\end{document}
